\documentclass[reqno]{amsart}
\usepackage[T1]{fontenc}
\usepackage{amsmath,amssymb,amsthm}
\usepackage{xr-hyper}
\usepackage[hidelinks]{hyperref}
\externaldocument[da-]{20261126T030300--the-data__workbook}
\externaldocument[ob-]{20261126T040100--spaces-forms-and-twists__workbook}
\externaldocument[dp-]{20261126T030900--denominators-of-projections-to-lifts__workbook}
\externaldocument[tw-]{20261126T031000--twists-with-few-congruences__workbook}
\externaldocument[w2-]{20261126T060600--lengths-as-values-of-polynomials__workbook}
\externaldocument[bk-]{20261126T080800--black-box-crystalline-lattices-and-the-condition-above-l__workbook}
\externaldocument[w4-]{20261126T081200--a-measure-for-hecke-values__workbook}

\newcounter{chap}
% Chapter numbers, generated from 20261125T012000--grh-workbook__guide.org
% by scripts/make-chapters.py.  Do not edit: change the order of the
% index and rerun the script (./preview-site.sh and CI do this).
%
% Usage in a sheet, after \newcounter{chap}:
%   % Chapter numbers, generated from 20261125T012000--grh-workbook__guide.org
% by scripts/make-chapters.py.  Do not edit: change the order of the
% index and rerun the script (./preview-site.sh and CI do this).
%
% Usage in a sheet, after \newcounter{chap}:
%   % Chapter numbers, generated from 20261125T012000--grh-workbook__guide.org
% by scripts/make-chapters.py.  Do not edit: change the order of the
% index and rerun the script (./preview-site.sh and CI do this).
%
% Usage in a sheet, after \newcounter{chap}:
%   \input{chapters}\chapnumber{<sheet ID>}
\makeatletter
\providecommand\chapters@set[2]{\@namedef{chapters@#1}{#2}}
\chapters@set{20261125T001000}{1}% zeros-to-ramp-sums
\chapters@set{20261125T001100}{2}% inert-primes-and-estimate-b
\chapters@set{20261125T005000}{3}% galois-fourier-and-averaging
\chapters@set{20261125T005500}{4}% columns-from-a-sphere
\chapters@set{20261125T000000}{5}% slope-method
\chapters@set{20261126T030100}{6}% a-twist-for-each-prime
\chapters@set{20261125T003000}{7}% the-chain-and-the-square-root
\chapters@set{20261125T010000}{8}% rank-lemma-and-zero-estimate
\chapters@set{20261125T010500}{9}% poincare-series-and-near-orthogonality
\chapters@set{20261124T223000}{10}% elliptic-curves-for-the-proof
\chapters@set{20261125T009500}{11}% fock-model-and-theta-intertwiners
\chapters@set{20261124T223100}{12}% complex-multiplication-and-inert-reduction
\chapters@set{20261125T008000}{13}% one-cm-determinant
\chapters@set{20261124T214500}{14}% supersingular-theta-divisibility
\chapters@set{20261124T214600}{15}% the-isotropic-plane
\chapters@set{20261125T013700}{16}% pairings-on-a-formal-group-of-height-two
\chapters@set{20261125T013500}{17}% hecke-nondegeneracy-on-the-tree
\chapters@set{20261126T030200}{18}% the-bound-is-attained
\chapters@set{20261125T009000}{19}% kronecker-limit-and-chowla-selberg
\chapters@set{20261125T002000}{20}% definite-theta-lifts
\chapters@set{20261125T015500}{21}% siegel-weil-in-the-smallest-case
\chapters@set{20261125T002100}{22}% definite-theta-lifts-and-norms
\chapters@set{20261125T016500}{23}% hecke-operators-and-theta-lifts
\chapters@set{20261125T016700}{24}% divisibility-along-the-chain
\chapters@set{20261125T004000}{25}% denominators-and-sha
\chapters@set{20261125T004100}{26}% tate-twists-and-sha
\chapters@set{20261126T020430}{27}% the-minus-part-by-gauss-sums
\chapters@set{20261125T014100}{28}% congruence-modules
\chapters@set{20261125T014500}{29}% eisenstein-congruences-and-selmer-classes
\chapters@set{20261125T014700}{30}% pseudocharacters-and-the-length-bound
\chapters@set{20261125T011000}{31}% galois-representations-black-box
\chapters@set{20261126T030300}{32}% the-data
\chapters@set{20261126T020210}{33}% bessel-poincare-series-and-their-gram-operators
\chapters@set{20261126T030400}{34}% the-seesaw-identity
\chapters@set{20261126T030500}{35}% divisibility-of-the-composite-lift
\chapters@set{20261126T030600}{36}% theorem-a-short-forms
\chapters@set{20261126T030700}{37}% the-determinant-of-a-chain
\chapters@set{20261126T030800}{38}% the-cost-of-one-lift
\chapters@set{20261126T030900}{39}% denominators-of-projections-to-lifts
\chapters@set{20261126T031000}{40}% twists-with-few-congruences
\chapters@set{20261126T031100}{41}% theorem-c-and-the-main-inequality
\providecommand\chapnumber[1]{%
  \@ifundefined{chapters@#1}%
    {\PackageError{chapters}{Sheet #1 is not listed in the index}%
      {List it in 20261125T012000--grh-workbook__guide.org\MessageBreak
       and run scripts/make-chapters.py.}}%
    {\edef\thechap{\@nameuse{chapters@#1}}}}
\makeatother
\chapnumber{<sheet ID>}
\makeatletter
\providecommand\chapters@set[2]{\@namedef{chapters@#1}{#2}}
\chapters@set{20261125T001000}{1}% zeros-to-ramp-sums
\chapters@set{20261125T001100}{2}% inert-primes-and-estimate-b
\chapters@set{20261125T005000}{3}% galois-fourier-and-averaging
\chapters@set{20261125T005500}{4}% columns-from-a-sphere
\chapters@set{20261125T000000}{5}% slope-method
\chapters@set{20261126T030100}{6}% a-twist-for-each-prime
\chapters@set{20261125T003000}{7}% the-chain-and-the-square-root
\chapters@set{20261125T010000}{8}% rank-lemma-and-zero-estimate
\chapters@set{20261125T010500}{9}% poincare-series-and-near-orthogonality
\chapters@set{20261124T223000}{10}% elliptic-curves-for-the-proof
\chapters@set{20261125T009500}{11}% fock-model-and-theta-intertwiners
\chapters@set{20261124T223100}{12}% complex-multiplication-and-inert-reduction
\chapters@set{20261125T008000}{13}% one-cm-determinant
\chapters@set{20261124T214500}{14}% supersingular-theta-divisibility
\chapters@set{20261124T214600}{15}% the-isotropic-plane
\chapters@set{20261125T013700}{16}% pairings-on-a-formal-group-of-height-two
\chapters@set{20261125T013500}{17}% hecke-nondegeneracy-on-the-tree
\chapters@set{20261126T030200}{18}% the-bound-is-attained
\chapters@set{20261125T009000}{19}% kronecker-limit-and-chowla-selberg
\chapters@set{20261125T002000}{20}% definite-theta-lifts
\chapters@set{20261125T015500}{21}% siegel-weil-in-the-smallest-case
\chapters@set{20261125T002100}{22}% definite-theta-lifts-and-norms
\chapters@set{20261125T016500}{23}% hecke-operators-and-theta-lifts
\chapters@set{20261125T016700}{24}% divisibility-along-the-chain
\chapters@set{20261125T004000}{25}% denominators-and-sha
\chapters@set{20261125T004100}{26}% tate-twists-and-sha
\chapters@set{20261126T020430}{27}% the-minus-part-by-gauss-sums
\chapters@set{20261125T014100}{28}% congruence-modules
\chapters@set{20261125T014500}{29}% eisenstein-congruences-and-selmer-classes
\chapters@set{20261125T014700}{30}% pseudocharacters-and-the-length-bound
\chapters@set{20261125T011000}{31}% galois-representations-black-box
\chapters@set{20261126T030300}{32}% the-data
\chapters@set{20261126T020210}{33}% bessel-poincare-series-and-their-gram-operators
\chapters@set{20261126T030400}{34}% the-seesaw-identity
\chapters@set{20261126T030500}{35}% divisibility-of-the-composite-lift
\chapters@set{20261126T030600}{36}% theorem-a-short-forms
\chapters@set{20261126T030700}{37}% the-determinant-of-a-chain
\chapters@set{20261126T030800}{38}% the-cost-of-one-lift
\chapters@set{20261126T030900}{39}% denominators-of-projections-to-lifts
\chapters@set{20261126T031000}{40}% twists-with-few-congruences
\chapters@set{20261126T031100}{41}% theorem-c-and-the-main-inequality
\providecommand\chapnumber[1]{%
  \@ifundefined{chapters@#1}%
    {\PackageError{chapters}{Sheet #1 is not listed in the index}%
      {List it in 20261125T012000--grh-workbook__guide.org\MessageBreak
       and run scripts/make-chapters.py.}}%
    {\edef\thechap{\@nameuse{chapters@#1}}}}
\makeatother
\chapnumber{<sheet ID>}
\makeatletter
\providecommand\chapters@set[2]{\@namedef{chapters@#1}{#2}}
\chapters@set{20261125T001000}{1}% zeros-to-ramp-sums
\chapters@set{20261125T001100}{2}% inert-primes-and-estimate-b
\chapters@set{20261125T005000}{3}% galois-fourier-and-averaging
\chapters@set{20261125T005500}{4}% columns-from-a-sphere
\chapters@set{20261125T000000}{5}% slope-method
\chapters@set{20261126T030100}{6}% a-twist-for-each-prime
\chapters@set{20261125T003000}{7}% the-chain-and-the-square-root
\chapters@set{20261125T010000}{8}% rank-lemma-and-zero-estimate
\chapters@set{20261125T010500}{9}% poincare-series-and-near-orthogonality
\chapters@set{20261124T223000}{10}% elliptic-curves-for-the-proof
\chapters@set{20261125T009500}{11}% fock-model-and-theta-intertwiners
\chapters@set{20261124T223100}{12}% complex-multiplication-and-inert-reduction
\chapters@set{20261125T008000}{13}% one-cm-determinant
\chapters@set{20261124T214500}{14}% supersingular-theta-divisibility
\chapters@set{20261124T214600}{15}% the-isotropic-plane
\chapters@set{20261125T013700}{16}% pairings-on-a-formal-group-of-height-two
\chapters@set{20261125T013500}{17}% hecke-nondegeneracy-on-the-tree
\chapters@set{20261126T030200}{18}% the-bound-is-attained
\chapters@set{20261125T009000}{19}% kronecker-limit-and-chowla-selberg
\chapters@set{20261125T002000}{20}% definite-theta-lifts
\chapters@set{20261125T015500}{21}% siegel-weil-in-the-smallest-case
\chapters@set{20261125T002100}{22}% definite-theta-lifts-and-norms
\chapters@set{20261125T016500}{23}% hecke-operators-and-theta-lifts
\chapters@set{20261125T016700}{24}% divisibility-along-the-chain
\chapters@set{20261125T004000}{25}% denominators-and-sha
\chapters@set{20261125T004100}{26}% tate-twists-and-sha
\chapters@set{20261126T020430}{27}% the-minus-part-by-gauss-sums
\chapters@set{20261125T014100}{28}% congruence-modules
\chapters@set{20261125T014500}{29}% eisenstein-congruences-and-selmer-classes
\chapters@set{20261125T014700}{30}% pseudocharacters-and-the-length-bound
\chapters@set{20261125T011000}{31}% galois-representations-black-box
\chapters@set{20261126T030300}{32}% the-data
\chapters@set{20261126T020210}{33}% bessel-poincare-series-and-their-gram-operators
\chapters@set{20261126T030400}{34}% the-seesaw-identity
\chapters@set{20261126T030500}{35}% divisibility-of-the-composite-lift
\chapters@set{20261126T030600}{36}% theorem-a-short-forms
\chapters@set{20261126T030700}{37}% the-determinant-of-a-chain
\chapters@set{20261126T030800}{38}% the-cost-of-one-lift
\chapters@set{20261126T030900}{39}% denominators-of-projections-to-lifts
\chapters@set{20261126T031000}{40}% twists-with-few-congruences
\chapters@set{20261126T031100}{41}% theorem-c-and-the-main-inequality
\providecommand\chapnumber[1]{%
  \@ifundefined{chapters@#1}%
    {\PackageError{chapters}{Sheet #1 is not listed in the index}%
      {List it in 20261125T012000--grh-workbook__guide.org\MessageBreak
       and run scripts/make-chapters.py.}}%
    {\edef\thechap{\@nameuse{chapters@#1}}}}
\makeatother
\chapnumber{20261126T081300}
\renewcommand{\thesection}{\thechap.\arabic{section}}

\theoremstyle{definition}
\newtheorem{problem}{Problem}[section]
\theoremstyle{plain}
\newtheorem{theorem}{Theorem}[section]
\theoremstyle{remark}
\newtheorem*{leanremark}{In the Lean formalization}

\newcommand{\Z}{\mathbb Z}
\newcommand{\Q}{\mathbb Q}
\newcommand{\F}{\mathbb F}
\newcommand{\C}{\mathbb C}
\newcommand{\Oc}{\mathcal O}
\newcommand{\TS}{T_{\mathrm{Sel}}}
\newcommand{\cyc}{\chi_{\mathrm{cyc}}}
\newcommand{\bd}{\mathrm{bd}}
\newcommand{\val}{\operatorname{val}}
\newcommand{\ord}{\operatorname{ord}}
\newcommand{\len}{\operatorname{len}}
\newcommand{\rk}{\operatorname{rk}}
\setlength{\emergencystretch}{3em}

\title{From a representation to a Selmer datum}
\author{}
\date{November 26, 2026}

\begin{document}
\maketitle
\refstepcounter{chap}\label{chap}

The small-twists theorem gives a single character whose first \(m\) powers bound the normalized Selmer lengths and the valuations of the Stickelberger system evaluations.  To use the small-twists theorem in the proof of Theorem~\ref{tw-thm:twisting}, we need to supply Selmer data for the representations and a bound for each datum.  We compute the global term in rank one, bound the three terms for a general lattice, then combine the Selmer data with the systems of the Hecke ratios.

\paragraph{Prerequisites.}
Use the data of Chapter~\ref{da-chap}, the chosen places, tame characters and critical pairs of Chapter~\ref{ob-chap}, and the main-summand-like condition, auxiliary twists and Selmer bound of Definitions~\ref{dp-def:main-summand-like}, \ref{dp-def:aux-twist} and~\ref{dp-def:selmer-bound}.  Use the Selmer datum, its three terms, characters for \(Q\) and torsion levels from Chapter~\ref{w2-chap}, and the crystalline cocycles and their rank from Chapter~\ref{bk-chap}.  We use Theorems~\ref{tw-thm:small} and~\ref{w4-thm:pair-system}.  For local reciprocity we use that local units map onto the abelianized inertia image, in either Frobenius convention.  The inverse avatar uses the convention of Definition~\ref{dp-def:aux-twist}.  We use elementary divisors over a discrete valuation ring \cite{La}.  Write \(\len_{\Oc}\) for ordinary module length and \(e_{\Oc}=\ord_{\Oc}(\ell)\), so normalized length is ordinary length divided by \(e_{\Oc}\).  Use \(G_E\), \(g^c\), the cyclotomic character \(\cyc\) and coefficient rings as in Section~\ref{dp-sec:statement}, with \(\ord_{\Oc}(0)=+\infty\).

\section{Rank-one calculations}\label{sec:rank-one}

The first three problems measure how long a character and its conjugate dual agree modulo powers of~\(3\).  For the rank-one calculation let \(E\) be any CM field, let \(\ell=3\), and let \(\Oc\) contain \(\Z_3[\mathrm i]\).  Give the rank-one lattice \(\Oc\) the action of a continuous character \(\chi\colon G_E\to\Oc^\times\).

\begin{problem}\label{pr:rank-one-dual}
For \(0\ne a\in\Oc\), use the conjugate additive dual of Definition~\ref{w2-def:selmer-datum}.  After choosing an \(\Oc\)-linear identification of the additive dual of \(\Oc/a\Oc\) with \(\Oc/a\Oc\), the equivariant maps to the conjugate dual correspond to the elements of \(\Oc/a\Oc\) killed by every difference
\[
  \chi(g)-\cyc(g^c)\chi((g^c)^{-1}),\qquad g\in G_E.
\]
\end{problem}

\begin{problem}\label{pr:rank-one-length}
Let \(j_0\ge 1\) be an integer and suppose \(\chi(g^c)\chi(g)=\cyc(g)^{-(j_0-1)}\) for every \(g\in G_E\).  The supremum, over \(0\ne a\in\Oc\), of the normalized lengths of the modules in Problem~\ref{pr:rank-one-dual} is \(\min_{g\in G_E}\val_3(\cyc(g)^{j_0}-1)\), with \(\val_3(0)=+\infty\).
\end{problem}

\begin{problem}\label{pr:rank-one-values}
For \(j_0=1,2,3,4,6\), show that the minimum of \(\val_3(u^{j_0}-1)\) over \(u\in\Z_3^\times\) is respectively \(0,1,0,1,2\).  At \(u=7\) the values are \(1,1,2,1,2\), and at \(u=19\) they are \(2,2,3,2,3\).
\end{problem}

Problem~\ref{pr:rank-one-values} computes the minimum in Problem~\ref{pr:rank-one-length} when \(\cyc(G_E)=\Z_3^\times\).  The cyclotomic minimum is of the same form as the number \(\nu_\ell\) in Zhang's account of the Selmer bound for finite-order twists \cite[\S5.2]{Zh}.

The next two problems measure the length of a module of rank-one cocycles over \(\Q_7\).  We use the classical presentation of the tame quotient by a generator \(s\) of tame inertia and an arithmetic Frobenius lift \(\mathrm{Fr}\), with relation \(\mathrm{Fr}\,s=s^q\,\mathrm{Fr}\), where \(q\) is the residue-field size \cite{SLF}.  In residue characteristic \(p\), wild inertia is pro-\(p\) and the finite quotients of tame inertia have order prime to \(p\).

\begin{problem}\label{pr:tame-model}
Over \(\Q_7\) take the rank-one module \(\Z/9\Z\), with wild inertia acting trivially, \(s\) acting by \(4\) and \(\mathrm{Fr}\) by \(u\in\{1,4,7\}\).  The pairs \((z(s),z(\mathrm{Fr}))\) of continuous cocycle values are exactly the solutions of
\[
  (1-4^7)z(\mathrm{Fr})=(1+4+\cdots+4^6-u)z(s)
  \quad\text{in }\Z/9\Z.
\]
\end{problem}

\begin{problem}\label{pr:tame-count}
For each \(u\) in Problem~\ref{pr:tame-model} there are \(27\) pairs, giving length \(3=2+\val_3(4-1)\) over \(\Z_3\).
\end{problem}

The next problem bounds the length of the restrictions of global cocycles to a decomposition group.

\begin{problem}\label{pr:local-bound}
For any prime \(\ell\), let \(L\) be an \(\Oc\)-lattice with continuous \(G_E\)-action.  Fix a finite place \(v\nmid\ell\) of residue characteristic \(p\) and a decomposition group \(G_v\) at \(v\).  Suppose some positive integer \(n_v\) prime to \(p\) satisfies \(L(h^{n_v})=1\) for all inertia elements \(h\) in \(G_v\), where \(L(g)\) denotes the matrix of the action.  Suppose also that \(L(h)=1\) for every inertia element \(h\) in some neighborhood of \(1\) in \(G_v\).  For \(0\ne a\in\Oc\) let \(C_v(a)\) be the module of restrictions to \(G_v\) of locally constant global cocycles with coefficients in \(L/aL\).  For every inertia element \(\delta\) in \(G_v\), prove
\[
  \begin{aligned}
    \len_{\Oc} C_v(a)&\le \len_{\Oc}(L/aL)+\len_{\Oc}\ker\bigl(L(\delta)-1\mid L/aL\bigr)\\
     &\le \len_{\Oc}(L/aL)+\ord_{\Oc}\det\bigl(L(\delta)-1\bigr).
  \end{aligned}
\]
\end{problem}

\begin{problem}\label{pr:scalar-model}
For multiplication by \(6\) on \((\Z/9\Z)^2\), the kernel has \(9\) elements and length \(2=2\val_3(6)\) over \(\Z_3\).
\end{problem}

\section{The three terms of a datum}\label{sec:bounds}

We now bound the global term, the sum of the local terms over the places dividing \(M\) prime to \(\ell\), and the two determinant valuations that make up the inertia term.

\begin{problem}\label{pr:glob}
Let \(L\) be a Galois lattice of positive rank \(r'\) whose action matrices span \(\operatorname{End}_{\Oc}(L)\).  Suppose \(g\) and \((g^c)^{-1}\) act by scalars \(\alpha\) and \(\alpha'\).  For each \(0\ne a\in\Oc\), on the equivariant maps \(L/aL\to (L/aL)^{*c}\) of Definition~\ref{w2-def:selmer-datum}, evaluation at a basis vector is injective, and every map is killed by \(\alpha-\cyc(g^c)\alpha'\).  Deduce that the supremum defining the global term is at most \(r'\ord_{\Oc}(\alpha-\cyc(g^c)\alpha')\).
\end{problem}

Fix \(\iota\colon\C\simeq\overline{\Q}_\ell\) and a main-summand-like \(R\) of rank \(r'\) for \((t',\mathrm{HT})\).  In the next problems take a lattice \(L\) of \(R\) over a coefficient ring large enough for clause~(e) of Definition~\ref{dp-def:main-summand-like}.  Take \(w_0\in\mathcal M\) with \(w_0\nmid\ell\), and take \(\Sigma_0\) to be the places dividing \(\ell MN\), with \(g_0=s_{w_0}\) the tame generator at the fixed place above \(w_0\).  Use the three terms of Definition~\ref{w2-def:selmer-datum} for \(L\) with \(\Sigma(\TS)\) the places dividing \(\ell M\).  As in Chapter~\ref{dp-chap}, \(n_\ell=\val_\ell(\mathcal N)\).

\begin{problem}\label{pr:global-budget}
If \(g\) acts by \(\alpha\) on \(L\), show that \((g^c)^{-1}\) acts by \(\cyc(g)^{\max(t'-1,0)}\alpha\).  For such a \(g\) with \(\cyc(g)=p_i^{r'}\) for some \(i\in\{0,1\}\), show that the scalar difference in Problem~\ref{pr:glob} is \(\alpha(1-p_i^{r'(\max(t'-1,0)+1)})\).  Prove that \(\mathrm{glob}/e_{\Oc}\le r'n_\ell\).
\end{problem}

\begin{problem}\label{pr:place-budget}
At most \(4f\) places divide \(M\), and \(\mathrm{plc}/e_{\Oc}\le 4fr'n_\ell\).
\end{problem}

\begin{problem}\label{pr:inertia-budget}
Show that the normalized determinant valuations at \(g_0\) and \((g_0^c)^{-1}\) in the definition of \(\mathrm{inr}\) are each at most \(r'n_\ell\).  Deduce \(\mathrm{inr}/e_{\Oc}\le 2r'n_\ell\) and \((\mathrm{glob}+\mathrm{plc}+\mathrm{inr})/e_{\Oc}\le(4f+3)r'n_\ell\).
\end{problem}

\section{The local conditions and the lattice}\label{sec:datum}

The remaining conditions for a Selmer datum are a scalar inertia action and a bound on the ranks of the crystalline cocycles.

\begin{problem}\label{pr:scalar}
Show that some inertia element at a place dividing \(M\) prime to \(\ell\) acts on \(L\) by a scalar \(\lambda\) with \(\lambda^2-1\in\Oc^\times\), and that its conjugate inverse acts by the same scalar.  To test the square-unit condition at \(\ell=3\), show that no unit whose residue belongs to \(\F_3^\times\) satisfies \(\lambda^2-1\in\Oc^\times\), and that \(\mathrm i\) satisfies the condition in \(\Z_3[\mathrm i]\).
\end{problem}

For a lattice in \(R\), the coefficient-ring form of the crystalline rank formula (Theorem~\ref{bk-thm:rank}) uses the labeled jumps of clause~(b).  At \(w\mid\ell\), the finitely generated \(\Oc\)-module \(\mathcal Z_w=Z_f(L|_{G_w})\) has rank
\[
  r'+\sum_{\sigma\mapsto w}\#\{u\in\mathrm{HT}_\sigma:u<0\},
\]
counting multiplicities.  Here \(\sigma:E\hookrightarrow\C\) runs over the embeddings inducing \(w\) through \(\iota\), and \(G_w\) is the fixed decomposition group at \(w\).  \emph{Status}: the labeled formula is read from the proof; not proved here.

\begin{problem}\label{pr:rank}
Show that \(\sum_{w\mid\ell}(\rk_{\Oc}\mathcal Z_w-r')\le fr'\).  With \(\mathcal Z_w\) as the modules of local cocycles, check the conditions of Definition~\ref{w2-def:selmer-datum} for \(L\).
\end{problem}

A continuous \(\ell\)-adic representation of \(G_E\) over \(\overline{\Q}_\ell\) has a lattice over every sufficiently large coefficient ring, meaning every coefficient ring containing one fixed ring.  \emph{Status}: read from the proof; not proved here.

Problems~\ref{pr:global-budget}--\ref{pr:rank}, with existence of lattices, prove the following.

\begin{theorem}\label{thm:datum}
Let \(R\) be main-summand-like of rank \(r'\) and \(w_0\in\mathcal M\) be prime to \(\ell\).  For all sufficiently large \(\Oc\), some lattice of \(R\) over \(\Oc\) carries a Selmer datum \(\TS\) for \((\Sigma_0,s_{w_0})\), where \(\Sigma_0\) consists of the places dividing \(\ell MN\).  The set \(\Sigma(\TS)\) consists of the places dividing \(\ell M\), the local relations are those of the crystalline cocycles, and \(\bd(\TS)\le (4f+3)r'n_\ell\).
\end{theorem}

\section{One character for all bounds}\label{sec:twists}

To apply the small-twists theorem at a common torsion level, choose \(w_0\) above one of \(p_0,p_1\) different from \(\ell\).  The set of chosen places above each \(p_i\) is nonempty and their residue-field sizes are \(p_i\) (Section~\ref{ob-sec:chosen}).  Use the same \(\Sigma_0\) and \(g_0\) for the representations and the systems of Theorem~\ref{w4-thm:pair-system}.  There are finitely many critical pairs (see the hint for Problem~\ref{pr:assemble}).

\begin{problem}\label{pr:torsion}
Given a finite family of main-summand-like representations, show that there is one coefficient ring over which each representation has a lattice with a datum as in Theorem~\ref{thm:datum}, and that \(c_{\mathrm{tor}}=(4f+3)rn_\ell+2\) is a torsion level for the family of data.  In the numerical model \(\ell=3\), \(r=e_{\Oc}=1\) and \(c_{\mathrm{tor}}=2\), the residue-field inequality asks for more than \(5\) elements.
\end{problem}

Two facts connect the character and length bounds of Theorem~\ref{tw-thm:small} to the twists used in Theorem~\ref{tw-thm:twisting}.

\begin{itemize}
\item \emph{The return from Galois characters to tame characters.}  For \(Q\) admissible for \(\Sigma_0\), every character for \(Q\) is the inverse avatar of a weight-zero tame character of modulus \(Q\).  \emph{Status}: read from the proof; not proved here.
\item \emph{Transport of crystalline relations.}  Let \(\Oc'\supseteq\Oc\) be coefficient rings and let a character over \(\Oc'\) be trivial on the fixed decomposition groups above \(\ell\).  Every crystalline relation on an \(\Oc\)-lattice, with coefficients mapped to \(\Oc'\), holds on the integral crystalline cocycles of the extended, twisted lattice.  \emph{Status}: read from the proof; not proved here.
\end{itemize}

The last two problems compare the two Selmer modules and account for the loss in passing from a system to a ratio.  Write \(\mathrm{tw}_\ell=(4f+3)rn_\ell+2^{-\ell}\) as in Definition~\ref{dp-def:clearing-bound}.

\begin{problem}\label{pr:assemble}
Let \(Q\) and \(\chi\) be given by Theorem~\ref{tw-thm:small} for \(\varepsilon=2^{-\ell}\), the data of Problem~\ref{pr:torsion} and the systems for all ramified critical pairs and signs.  Show that \(\chi\) is the inverse avatar of a tame character \(\vartheta\) of modulus \(Q\), and that for \(1\le j\le m\), \(\vartheta^j\) ramifies at every \(w\in Q\) and has inverse avatar \(\chi^j\).  For each datum, each such \(j\) and each sufficiently large coefficient ring \(\Oc'\), take \(L_{\mathrm{tw}}\) to be the datum's lattice extended to \(\Oc'\) and twisted by \(\chi^j\).  Show that \(\mathrm{Sel}(L_{\mathrm{tw}},\gamma(\TS))\) of Definition~\ref{dp-def:selmer-bound} is a submodule of the module of Definition~\ref{w2-def:twisted-length} for the same datum and twist.  Deduce clause~2 of Theorem~\ref{tw-thm:twisting}.
\end{problem}

\begin{problem}\label{pr:twisting}
Prove \(4\val_\ell(\#k_{w_0}-1)\le(4f+3)rn_\ell\).  Show that \(\vartheta\) from Problem~\ref{pr:assemble} is an \(m\)-fold twist as in Section~\ref{ob-sec:twists} and that each power \(\vartheta^j\) for \(1\le j\le m\) is auxiliary with ramification at most \(2^{-\ell}\).  Show that each ratio valuation in clause~3 of Theorem~\ref{tw-thm:twisting} is at most \(4\val_\ell(\#k_{w_0}-1)+2^{-\ell}\).  Deduce clauses~1 and~3 of that theorem.
\end{problem}

\paragraph{Taken as given.}
We used the tame presentation, elementary divisors, the labeled crystalline rank formula, existence of lattices, the return from Galois characters to tame characters and transport of crystalline relations.  The small-twists theorem and the systems of the ratios supply the two inputs to Problems~\ref{pr:assemble} and~\ref{pr:twisting}, which finish the deduction of Theorem~\ref{tw-thm:twisting}.

\begin{leanremark}
In \cite{GL}, Theorem~\ref{thm:datum} is
\nolinkurl{IsMainSummandLike.eventually_exists_selmerDatum}.
The global bound uses \nolinkurl{congruenceTerm_le} and
\nolinkurl{IsMainSummandLike.conjGal_inv_eq_smul_one}.
Lattice existence is \nolinkurl{eventually_exists_isLatticeIn} and the labeled rank formula is
\nolinkurl{finrank_crysCocycles_eq_add_neg_jumps}.
The assembly uses \nolinkurl{exists_selmerData_isTorsionLevel},
\nolinkurl{exists_tameChar_invAvatar_eq},
\nolinkurl{mapRel_crysRelations_subset} in
\nolinkurl{exists_mFoldTwist_hasDivisibleSelmerLengthLe}.
See also \nolinkurl{exists_stickelbergerSystem_ellVal_ratio_le}.
\end{leanremark}

\section*{Hints}

Use the hints for the steps involving the tame relation, the trace pairing and the comparison after twisting.

\begin{description}
\item[\ref{pr:tame-model}] Express the cocycle condition as a homomorphism to the group of affine maps over \(\Z/9\Z\).
\item[\ref{pr:local-bound}] Use the tame presentation and elementary divisors.
\item[\ref{pr:glob}] Use the matrix units and the additivity of length.
\item[\ref{pr:global-budget}] Test the scalar identity against \(L(h)\) in the trace pairing.  For the valuation use the order of the general linear group in Section~\ref{ob-sec:chosen}.
\item[\ref{pr:scalar}] Use clause~(f) and the cyclotomic character on inertia away from \(\ell\).
\item[\ref{pr:torsion}] Compare the two inequalities in Definition~\ref{w2-def:torsion-level} as the residue degree grows.
\item[\ref{pr:assemble}] For finiteness, fix the weight, tame part and \(z\).  Use \(y\mapsto y/\bar y\) from the ideal class group to \(\mathrm{Cl}_\varnothing\) of Section~\ref{w2-sec:systems}.  For the local conditions use the Artin formula on local units and on ideles supported above \(\ell\).
\item[\ref{pr:twisting}] Use the implication of Problem~\ref{w2-pr:ratio-from-system}.  For the ramification bound, restrict the sum in the small-twists theorem to \(Q\).
\end{description}

\begin{thebibliography}{SLF}
\bibitem[GL]{GL} L.~Sernau, \emph{GRH for Dirichlet \(L\)-functions, in Lean~4}, a shorter rewriting of the Lean proof, text by Claude, \url{https://github.com/lsernau/grh-lean}, commit \texttt{72fead1}, readable edition \texttt{grh/}.
\bibitem[La]{La} S.~Lang, \emph{Algebra}, 3rd ed., Grad.\ Texts in Math.\ 211, Springer, 2002.
\bibitem[SLF]{SLF} J.-P.~Serre, \emph{Local Fields}, Grad.\ Texts in Math.\ 67, Springer, 1979.
\bibitem[Zh]{Zh} R.~Zhang, \emph{A preliminary note on Alp\"oge's proposed proof of the generalized Riemann hypothesis}, Hexagon v2, \url{https://hexagonmath.org/pdf/2611.00809v2}.
\end{thebibliography}

\end{document}
